% ─────────────────────────────────────────────────────────────────────────────
% Statuten eines Vereins "Erneuerbare-Energie-Gemeinschaft" (Statuten-Generator, webapp:
% LatexStatutePdfRenderer, fachlich: docs/STATUTEN_GENERATOR.md).
%
% Regeln für diese Vorlage:
%  - Querverweise AUSSCHLIESSLICH mit \label/\ref (nie händische Paragraphen-Nummern).
%    \ref auf einen Paragraphen ergibt "§ n", auf einen Absatz "§ n Abs m".
%  - RAW_-Variablen kommen fertig escaped aus PHP (wahlabhängige Absätze, inkl. \item/\label).
%  - Kopfzeile "Entwurf – keine Rechtsberatung" (RAW_ENTWURF) bis zur Freigabe durch den Admin.
%  - Rechtsstand: VerG 2002, ElWG (BGBl. I 91/2025), EAG § 79 -- vor Freigabe im RIS prüfen.
% ─────────────────────────────────────────────────────────────────────────────
\documentclass[a4paper,11pt]{article}
\usepackage[utf8]{inputenc}
\usepackage[T1]{fontenc}
\usepackage[ngerman]{babel}
\usepackage[left=25mm,right=25mm,top=25mm,bottom=25mm]{geometry}
\usepackage{microtype}
\usepackage{parskip}
\usepackage{enumitem}
\usepackage{fancyhdr}
\usepackage{lastpage}
\usepackage{xcolor}
\usepackage[scaled=0.96]{helvet}
\renewcommand{\familydefault}{\sfdefault}
\usepackage{titlesec}
\usepackage[hidelinks]{hyperref}

\definecolor{navy}{RGB}{26,32,51}
\definecolor{midgray}{RGB}{107,114,128}
\definecolor{warnred}{RGB}{185,28,28}

% Paragraphen: "§ 1 Name, Sitz ..." -- \ref liefert "§ 1"
\renewcommand{\thesection}{\S\,\arabic{section}}
\titleformat{\section}{\large\bfseries\color{navy}}{\thesection}{0.8em}{}
\titlespacing*{\section}{0pt}{5mm}{1.5mm}
% Absätze: "(1)" -- \ref auf einen Absatz liefert "§ n Abs m"
\newlist{absaetze}{enumerate}{1}
\setlist[absaetze]{label=(\arabic*),ref=\thesection~Abs~\arabic*,leftmargin=9mm,itemsep=1mm,topsep=1mm}

\def\entwurf{<<<RAW_ENTWURF>>>}

\pagestyle{fancy}
\fancyhf{}
\renewcommand{\headrulewidth}{0pt}
\ifx\entwurf\empty\else
  \fancyhead[C]{\small\bfseries\color{warnred}\entwurf}
\fi
\fancyfoot[L]{\scriptsize\color{midgray}Statuten <<<VEREINSNAME>>> -- Fassung <<<VERSION>>> vom <<<DATUM>>>}
\fancyfoot[R]{\scriptsize\color{midgray}Seite \thepage{} von \pageref{LastPage}}

% Deckblatt ohne Fußzeile; die Entwurfs-Kopfzeile bleibt bis zur Freigabe sichtbar.
\fancypagestyle{deckblatt}{%
  \fancyhf{}%
  \ifx\entwurf\empty\else\fancyhead[C]{\small\bfseries\color{warnred}\entwurf}\fi
}
% Inhaltsverzeichnis: Platz für "§ 10" usw.
\makeatletter
\renewcommand*\l@section{\@dottedtocline{1}{0em}{2.8em}}
\makeatother

\begin{document}

\thispagestyle{deckblatt}
\begin{center}
  \vspace*{3cm}
  {\Huge\bfseries\color{navy}STATUTEN\par}
  \vspace{0.5cm}
  {\Large des Vereins\par}
  \vspace{1cm}
  {\LARGE\bfseries Erneuerbare-Energie-Gemeinschaft\par}
  \vspace{0.3cm}
  {\Large\bfseries <<<DECKBLATT_NAME>>>\par}
  \vspace{2cm}
  {\large <<<NAHEBEREICH_ART>>> Erneuerbare-Energie-Gemeinschaft gemäß § 79 EAG und §§ 66 ff ElWG\par}
  \vspace{0.5cm}
  <<<RAW_DECKBLATT_ID>>>
  \vspace{0.5cm}
  {\large Netzbetreiber: <<<NETZBETREIBER>>>\par}
\end{center}

\newpage
\tableofcontents
\newpage

% ─── § 1 ───────────────────────────────────────────────────────────────────
\section{Name, Sitz und Tätigkeitsbereich}\label{par:name}
\begin{absaetze}
  \item Der Verein führt den Namen \glqq{}<<<VEREINSNAME>>>\grqq{}.
  \item Er hat seinen Sitz in <<<SITZ>>> und erstreckt seine Tätigkeit auf <<<RAW_TAETIGKEITSBEREICH>>>.
  \item Der Verein ist eine Erneuerbare-Energie-Gemeinschaft im Sinne des § 79 EAG und des § 66 ElWG. Die Anlagen
        der Mitglieder befinden sich im <<<NAHEBEREICH>>> (§ 70 Abs 6 ElWG) des Verteilernetzes der <<<NETZBETREIBER>>><<<RAW_NAHEBEREICH_ID>>>.\label{abs:nahebereich}
  \item Die Errichtung von Zweigvereinen ist nicht beabsichtigt.
\end{absaetze}

% ─── § 2 ───────────────────────────────────────────────────────────────────
\section{Zweck}\label{par:zweck}
\begin{absaetze}
  \item Der Verein bezweckt die gemeinsame Erzeugung, Speicherung, den Verbrauch und den Verkauf von Energie aus
        erneuerbaren Quellen (<<<ENERGIETRAEGER>>>) durch und für seine Mitglieder.
  \item Hauptzweck ist nicht die Erzielung finanzieller Gewinne, sondern es sollen den Mitgliedern und dem Gebiet, in
        dem der Verein tätig ist, ökologische, wirtschaftliche und sozialgemeinschaftliche Vorteile verschafft werden.
  \item Der Verein ist nicht gemeinnützig im Sinne der §§ 34 ff BAO.\label{abs:nichtgemeinnuetzig}
  \item Der Verein dient der Energiewende, der Stärkung der regionalen Versorgung und der Bewusstseinsbildung für einen
        sparsamen und effizienten Umgang mit Energie.
\end{absaetze}

% ─── § 3 ───────────────────────────────────────────────────────────────────
\section{Mittel zur Erreichung des Vereinszwecks}\label{par:mittel}
\begin{absaetze}
  \item Der Vereinszweck soll durch die in \ref{abs:ideell} und \ref{abs:materiell} angeführten ideellen und
        materiellen Mittel erreicht werden.
  \item Als ideelle Mittel dienen\label{abs:ideell}
        \begin{enumerate}[label=\alph*)]
          \item die Organisation der gemeinsamen Nutzung der von den Mitgliedern oder vom Verein erzeugten Energie
                (<<<ANLAGENMODELL>>>),
          \item die Abwicklung des Datenaustauschs mit dem Netzbetreiber und die Abrechnung unter den Mitgliedern,
          \item Information, Beratung und Veranstaltungen.
        \end{enumerate}
  \item Die erforderlichen materiellen Mittel werden aufgebracht durch\label{abs:materiell}
        \begin{enumerate}[label=\alph*)]
          \item Entgelte für die von den Mitgliedern bezogene Energie,
          \item Mitgliedsbeiträge<<<RAW_EINTRITTSGELD>>>,
          \item Förderungen, Spenden, Zuwendungen und sonstige Erträge.
        \end{enumerate}
  \item <<<RAW_PREISE>>>
  \item Etwaige Überschüsse dürfen nur für den Vereinszweck verwendet werden, insbesondere zur Senkung der Entgelte.
        Eine Ausschüttung von Gewinnen an die Mitglieder ist ausgeschlossen. Die Mitglieder haben bei ihrem Ausscheiden
        und bei Auflösung des Vereins keinen Anspruch auf das Vereinsvermögen.
  \item <<<RAW_NACHSCHUSS>>>
\end{absaetze}

% ─── § 4 ───────────────────────────────────────────────────────────────────
\section{Arten der Mitgliedschaft}\label{par:arten}
\begin{absaetze}
  \item <<<RAW_MITGLIEDSARTEN>>>
  \item Ordentliche Mitglieder nehmen mit mindestens einer Verbrauchs- oder Erzeugungsanlage (Zählpunkt) an der
        Erneuerbare-Energie-Gemeinschaft teil.
\end{absaetze}

% ─── § 5 ───────────────────────────────────────────────────────────────────
\section{Erwerb der Mitgliedschaft}\label{par:erwerb}
\begin{absaetze}
  \item Ordentliche Mitglieder können natürliche Personen, Gemeinden, Rechtsträger von Behörden in Bezug auf lokale
        Dienststellen, sonstige juristische Personen des öffentlichen Rechts sowie kleine und mittlere Unternehmen sein
        (§ 6 Abs 1 Z 45 ElWG), deren Anlagen sich im Nahebereich gemäß \ref{abs:nahebereich} befinden. Große Unternehmen
        können nicht ordentliche Mitglieder sein. Bei privaten Unternehmen darf die Teilnahme nicht ihre gewerbliche
        oder berufliche Haupttätigkeit sein.\label{abs:mitgliedkreis}
  \item Haushalte, die als schutzbedürftig gelten oder von Energiearmut betroffen sind, dürfen von der Mitgliedschaft
        nicht ausgeschlossen oder benachteiligt werden.\label{abs:schutzbeduerftig}
  \item <<<RAW_AUFNAHME>>>
  \item Die Mitgliedschaft wird mit der Aufnahme wirksam; die Teilnahme an der gemeinsamen Energienutzung beginnt mit
        der Zuordnung des Zählpunkts durch den Netzbetreiber.
  \item Bis zur Entstehung des Vereins erfolgt die vorläufige Aufnahme durch die Gründungsmitglieder.
\end{absaetze}

% ─── § 6 ───────────────────────────────────────────────────────────────────
\section{Beendigung der Mitgliedschaft}\label{par:ende}
\begin{absaetze}
  \item Die Mitgliedschaft erlischt durch Tod, bei juristischen Personen durch Verlust der Rechtspersönlichkeit,
        durch freiwilligen Austritt, durch Ausschluss und durch Wegfall der Voraussetzungen nach
        \ref{abs:mitgliedkreis}.
  \item <<<RAW_RECHTSNACHFOLGE>>>
  \item Der Austritt ist jederzeit <<<AUSTRITTSFRIST>>> möglich und dem Vorstand schriftlich (auch per E-Mail)
        mitzuteilen. Zwingende gesetzliche Kündigungs- und Rücktrittsrechte, insbesondere für Verbraucher und
        Kleinunternehmen, bleiben unberührt und gehen dieser Bestimmung vor.\label{abs:austritt}
  \item <<<RAW_AUSSCHLUSS>>>
  \item Gegen den Ausschluss kann das Mitglied binnen vier Wochen die Schlichtungseinrichtung (\ref{par:schlichtung})
        anrufen.
  \item Mit dem Ende der Mitgliedschaft endet die Teilnahme an der gemeinsamen Energienutzung. Bis dahin entstandene
        Forderungen und Verbindlichkeiten werden in der nächsten Abrechnung ausgeglichen.
\end{absaetze}

% ─── § 7 ───────────────────────────────────────────────────────────────────
\section{Rechte und Pflichten der Mitglieder}\label{par:rechte}
\begin{absaetze}
  \item Die Mitglieder sind berechtigt, an allen Veranstaltungen des Vereins teilzunehmen und die Einrichtungen des
        Vereins zu beanspruchen. Das Stimmrecht in der Mitgliederversammlung sowie das aktive und passive Wahlrecht
        stehen nur den ordentlichen Mitgliedern zu.
  \item Jedes Mitglied kann vom Vorstand die Ausfolgung der Statuten verlangen. Mindestens ein Zehntel der Mitglieder
        kann vom Vorstand die Einberufung einer Mitgliederversammlung verlangen (\ref{abs:aomv}).
  \item Die Mitglieder sind in jeder Mitgliederversammlung vom Vorstand über die Tätigkeit und die finanzielle
        Gebarung des Vereins zu informieren. Der geprüfte Rechnungsabschluss ist den Mitgliedern auf Verlangen
        zugänglich zu machen.
  \item Die Mitglieder sind verpflichtet, die Interessen des Vereins zu fördern, die Statuten und die Beschlüsse der
        Vereinsorgane zu beachten und die Mitgliedsbeiträge sowie die Entgelte pünktlich zu bezahlen.
  \item <<<RAW_DATENFREIGABE>>>
  \item Der Verein rechnet die Energie spätestens <<<ABRECHNUNGSFRIST>>> Tage nach Ende der Abrechnungsperiode ab.
\end{absaetze}

% ─── § 8 ───────────────────────────────────────────────────────────────────
\section{Vereinsorgane}\label{par:organe}
Organe des Vereins sind die Mitgliederversammlung (\ref{par:mv}, \ref{par:mvaufgaben}), der Vorstand
(\ref{par:vorstand} bis \ref{par:obliegenheiten}), die Rechnungsprüfer (\ref{par:rp}) und die Schlichtungseinrichtung
(\ref{par:schlichtung}).

% ─── § 9 ───────────────────────────────────────────────────────────────────
\section{Mitgliederversammlung}\label{par:mv}
\begin{absaetze}
  \item Die Mitgliederversammlung ist die \glqq{}Mitgliederversammlung\grqq{} im Sinne des Vereinsgesetzes 2002. Eine
        ordentliche Mitgliederversammlung findet <<<MV_INTERVALL>>> statt.
  \item Eine außerordentliche Mitgliederversammlung findet auf Beschluss des Vorstands oder der ordentlichen
        Mitgliederversammlung, auf schriftlichen Antrag von mindestens einem Zehntel der Mitglieder oder auf
        Verlangen der Rechnungsprüfer binnen <<<AOMV_WOCHEN>>> Wochen statt.\label{abs:aomv}
  \item Zu den Mitgliederversammlungen sind alle Mitglieder mindestens <<<EINLADUNG_TAGE>>> Tage vor dem Termin
        <<<RAW_EINLADUNG>>> unter Angabe der Tagesordnung einzuladen. Die Einberufung erfolgt durch den Vorstand.
  \item Anträge zur Mitgliederversammlung sind mindestens drei Tage vor dem Termin beim Vorstand schriftlich
        einzureichen. Gültige Beschlüsse, ausgenommen über einen Antrag auf Einberufung einer außerordentlichen
        Mitgliederversammlung, können nur zur Tagesordnung gefasst werden.
  \item <<<RAW_VIRTUELL>>>
  \item <<<RAW_STIMMRECHT>>>
  \item <<<RAW_BESCHLUSSFAEHIGKEIT>>>
  \item Die Wahlen und die Beschlussfassungen erfolgen in der Regel mit einfacher Mehrheit der abgegebenen gültigen
        Stimmen. Beschlüsse, mit denen die Statuten geändert oder der Verein aufgelöst werden soll, bedürfen einer
        qualifizierten Mehrheit von <<<MEHRHEIT>>> der abgegebenen gültigen Stimmen.\label{abs:mehrheit}
  \item Den Vorsitz führt der Obmann bzw. die Obfrau, bei Verhinderung die Stellvertretung nach
        \ref{abs:vertretung}. Ist auch diese verhindert, führt das an Jahren älteste anwesende Vorstandsmitglied den
        Vorsitz.
\end{absaetze}

% ─── § 10 ──────────────────────────────────────────────────────────────────
\section{Aufgaben der Mitgliederversammlung}\label{par:mvaufgaben}
Der Mitgliederversammlung sind folgende Aufgaben vorbehalten:
\begin{enumerate}[label=\alph*)]
  \item Entgegennahme und Genehmigung des Rechenschaftsberichts und des Rechnungsabschlusses unter Einbindung der
        Rechnungsprüfer,
  \item Wahl und Enthebung der Mitglieder des Vorstands und der Rechnungsprüfer,
  \item Genehmigung von Rechtsgeschäften zwischen Rechnungsprüfern und Verein,
  \item Entlastung des Vorstands,
  \item <<<RAW_MV_ENTGELTE>>>
<<<RAW_MV_EHREN>>>
  \item Beschlussfassung über Statutenänderungen und über die freiwillige Auflösung des Vereins
        (\ref{abs:mehrheit}),
  \item Beratung und Beschlussfassung über sonstige auf der Tagesordnung stehende Fragen.
\end{enumerate}

% ─── § 11 ──────────────────────────────────────────────────────────────────
\section{Vorstand}\label{par:vorstand}
\begin{absaetze}
  \item <<<RAW_VORSTAND_BESETZUNG>>>
  \item Der Vorstand wird von der Mitgliederversammlung gewählt. Die Funktionsperiode beträgt <<<FUNKTIONSPERIODE>>>;
        sie dauert jedenfalls bis zur Wahl eines neuen Vorstands. <<<RAW_WIEDERWAHL>>>
  \item <<<RAW_KOOPTIERUNG>>>
  \item Der Vorstand wird vom Obmann bzw. der Obfrau, bei Verhinderung von der Stellvertretung nach
        \ref{abs:vertretung}, schriftlich oder mündlich einberufen.
  \item Der Vorstand ist beschlussfähig, wenn alle seine Mitglieder eingeladen wurden und mindestens die Hälfte von
        ihnen anwesend ist. <<<RAW_VORSTAND_BESCHLUSS>>>
  \item <<<RAW_UMLAUF>>>
  \item Die Funktion eines Vorstandsmitglieds erlischt durch Tod, Ablauf der Funktionsperiode, Enthebung durch die
        Mitgliederversammlung und Rücktritt. Der Rücktritt ist dem Vorstand, im Fall des Rücktritts des gesamten
        Vorstands der Mitgliederversammlung schriftlich zu erklären.
  \item Die Vorstandsmitglieder üben ihre Funktion ehrenamtlich aus; der Ersatz angemessener Auslagen ist zulässig.
\end{absaetze}

% ─── § 12 ──────────────────────────────────────────────────────────────────
\section{Aufgaben des Vorstands}\label{par:vorstandaufgaben}
Dem Vorstand obliegt die Leitung des Vereins. Er ist das \glqq{}Leitungsorgan\grqq{} im Sinne des Vereinsgesetzes 2002.
Ihm kommen alle Aufgaben zu, die nicht durch die Statuten einem anderen Vereinsorgan zugewiesen sind, insbesondere:
\begin{enumerate}[label=\alph*)]
  \item Einrichtung eines den Anforderungen des Vereins entsprechenden Rechnungswesens mit laufender Aufzeichnung der
        Einnahmen und Ausgaben und Führung eines Vermögensverzeichnisses,
  \item Erstellung des Jahresvoranschlags, des Rechenschaftsberichts und des Rechnungsabschlusses binnen fünf Monaten
        nach Ende des Rechnungsjahres (§ 21 VerG),
  \item Vorbereitung und Einberufung der Mitgliederversammlung (\ref{par:mv}),
  \item <<<RAW_VORSTAND_ENTGELTE>>>
  \item Aufnahme und Ausschluss von Mitgliedern, soweit nicht der Mitgliederversammlung vorbehalten
        (\ref{par:erwerb}, \ref{par:ende}),
  \item Abschluss der Verträge mit dem Netzbetreiber und den Mitgliedern sowie die Abwicklung der gemeinsamen
        Energienutzung.
\end{enumerate}

% ─── § 13 ──────────────────────────────────────────────────────────────────
\section{Besondere Obliegenheiten einzelner Vorstandsmitglieder}\label{par:obliegenheiten}
\begin{absaetze}
  \item Der Obmann bzw. die Obfrau führt die laufenden Geschäfte des Vereins. Der Kassier bzw. die Kassierin ist für
        die ordnungsgemäße Geldgebarung verantwortlich.
  \item <<<RAW_VERTRETUNG_AUSSEN>>>
  \item Rechtsgeschäfte über einem Betrag von <<<VIER_AUGEN>>> Euro bedürfen der gemeinsamen Zeichnung von Obmann bzw.
        Obfrau und Kassier bzw. Kassierin. Rechtsgeschäfte zwischen Vorstandsmitgliedern und dem Verein bedürfen der
        Zustimmung eines anderen, zur Vertretung befugten Vorstandsmitglieds (§ 6 Abs 4 VerG).\label{abs:vieraugen}
  \item Bei Gefahr im Verzug kann der Obmann bzw. die Obfrau auch in Angelegenheiten, die in den Wirkungsbereich der
        Mitgliederversammlung oder des Vorstands fallen, selbständig Anordnungen treffen; diese bedürfen der
        nachträglichen Genehmigung durch das zuständige Vereinsorgan.
  \item <<<RAW_SCHRIFTFUEHRER>>>
  \item <<<RAW_STELLVERTRETUNG>>>\label{abs:vertretung}
\end{absaetze}

% ─── § 14 ──────────────────────────────────────────────────────────────────
\section{Rechnungsprüfer}\label{par:rp}
\begin{absaetze}
  \item <<<RP_ANZAHL>>> Rechnungsprüfer werden von der Mitgliederversammlung für die Dauer von <<<RP_PERIODE>>>
        gewählt. Wiederwahl ist möglich. Die Rechnungsprüfer müssen nicht Mitglied des Vereins sein.
  \item Die Rechnungsprüfer müssen unabhängig und unbefangen sein und dürfen keinem Organ – mit Ausnahme der
        Mitgliederversammlung – angehören, dessen Tätigkeit Gegenstand der Prüfung ist (§ 5 Abs 5 VerG).
  \item Den Rechnungsprüfern obliegen die laufende Geschäftskontrolle sowie die Prüfung der Finanzgebarung des Vereins
        im Hinblick auf die Ordnungsmäßigkeit der Rechnungslegung und die statutengemäße Verwendung der Mittel binnen
        vier Monaten ab Erstellung der Einnahmen- und Ausgabenrechnung. Sie berichten der Mitgliederversammlung über
        das Ergebnis der Prüfung.
  \item Überschreiten die gewöhnlichen Einnahmen oder Ausgaben die Schwellenwerte des § 22 VerG, gelten die dort
        vorgesehenen Pflichten (Jahresabschluss bzw. Abschlussprüfer) zusätzlich.
\end{absaetze}

% ─── § 15 ──────────────────────────────────────────────────────────────────
\section{Schlichtungseinrichtung}\label{par:schlichtung}
\begin{absaetze}
  \item Zur Schlichtung von allen aus dem Vereinsverhältnis entstehenden Streitigkeiten ist die vereinsinterne
        Schlichtungseinrichtung berufen. Sie ist eine \glqq{}Schlichtungseinrichtung\grqq{} im Sinne des § 8 VerG und kein
        Schiedsgericht nach den §§ 577 ff ZPO.
  \item Die Schlichtungseinrichtung setzt sich aus drei ordentlichen Mitgliedern zusammen. Jeder Streitteil macht
        binnen sieben Tagen ein Mitglied namhaft; diese wählen binnen weiterer 14 Tage ein drittes Mitglied zum
        Vorsitzenden. Die Mitglieder der Schlichtungseinrichtung dürfen keinem Organ – mit Ausnahme der
        Mitgliederversammlung – angehören, dessen Tätigkeit Gegenstand der Streitigkeit ist.
  \item Die Schlichtungseinrichtung entscheidet nach Gewährung beiderseitigen Gehörs bei Anwesenheit aller ihrer
        Mitglieder mit einfacher Mehrheit. Ihre Entscheidungen sind vereinsintern endgültig.
  \item Der ordentliche Rechtsweg steht nach Ablauf von sechs Monaten ab Anrufung der Schlichtungseinrichtung offen,
        sofern das Verfahren nicht früher beendet ist (§ 8 Abs 1 VerG).
\end{absaetze}

% ─── § 16 ──────────────────────────────────────────────────────────────────
\section{Datenschutz}\label{par:datenschutz}
\begin{absaetze}
  \item Der Verein verarbeitet personenbezogene Daten seiner Mitglieder (insbesondere Name, Anschrift,
        Kontaktdaten, Bankverbindung, Zählpunktnummern und Energiedaten) im Einklang mit der DSGVO und dem DSG,
        soweit dies für die Mitgliederverwaltung, die gemeinsame Energienutzung und die Abrechnung erforderlich ist.
  \item Die Daten werden an den Netzbetreiber und über die Energiewirtschaftliche Datenaustauschplattform (EDA)
        übermittelt, soweit dies gesetzlich oder für die Teilnahme erforderlich ist. Mit Dienstleistern, die Daten im
        Auftrag des Vereins verarbeiten, werden Auftragsverarbeitungsverträge geschlossen (Art 28 DSGVO).
  \item Buchhaltungsunterlagen werden für die Dauer der gesetzlichen Aufbewahrungspflicht (§ 132 BAO, sieben Jahre)
        aufbewahrt; sonstige Daten werden nach dem Ende der Mitgliedschaft gelöscht, sobald sie nicht mehr benötigt
        werden.
\end{absaetze}

% ─── § 17 ──────────────────────────────────────────────────────────────────
\section{Freiwillige Auflösung des Vereins}\label{par:aufloesung}
\begin{absaetze}
  \item Die freiwillige Auflösung des Vereins kann nur in einer zu diesem Zweck einberufenen Mitgliederversammlung
        mit der Mehrheit nach \ref{abs:mehrheit} beschlossen werden.
  \item Diese Mitgliederversammlung hat – sofern Vereinsvermögen vorhanden ist – über die Abwicklung zu beschließen,
        insbesondere einen Abwickler zu berufen.
  \item Das nach Abdeckung der Verbindlichkeiten verbleibende Vereinsvermögen ist <<<AUFLOESUNG_VERMOEGEN>>> zu
        übertragen bzw. zu verwenden. Eine Verteilung an die Mitglieder ist ausgeschlossen (\ref{par:mittel}).
  \item Der letzte Vereinsvorstand hat die freiwillige Auflösung binnen vier Wochen nach Beschlussfassung der
        zuständigen Vereinsbehörde schriftlich anzuzeigen (§ 28 VerG). Die Abmeldung beim Netzbetreiber ist rechtzeitig
        zu veranlassen.
\end{absaetze}

% ─── § 18 ──────────────────────────────────────────────────────────────────
\section{Schlussbestimmungen}\label{par:schluss}
\begin{absaetze}
  \item Personenbezogene Bezeichnungen gelten für alle Geschlechter gleichermaßen.
  \item Sollten einzelne Bestimmungen dieser Statuten unwirksam sein, bleibt die Wirksamkeit der übrigen Bestimmungen
        unberührt. Anstelle der unwirksamen Bestimmung gilt die gesetzliche Regelung.
\end{absaetze}

\vspace{8mm}
{\small\color{midgray}<<<RAW_HINWEIS>>>}

\end{document}
